% ============================================================================
% BAB IV - PENGUJIAN DAN ANALISIS
% ============================================================================
\chapter{PENGUJIAN DAN ANALISIS}

Bab ini menyajikan hasil pengujian sistem yang difokuskan pada dua
skenario utama: (A) pengujian mekanisme \textit{read-write splitting}
untuk membuktikan bahwa query \textit{write} dan \textit{read}
diproses oleh \textit{server} yang berbeda, dan (B) pengujian
toleransi kesalahan (\textit{fault tolerance}) dengan mematikan
salah satu Replica secara sengaja untuk memastikan aplikasi tetap
berjalan. Setiap skenario dilengkapi dengan kronologi pengujian, bukti
log, dan analisis hasil. Kedua skenario ini dipilih karena mewakili
dua aspek paling kritis dari arsitektur \textit{read-write splitting}:
kebenaran fungsional (apakah query diarahkan ke node yang tepat) dan
ketahanan sistem (apakah sistem tetap berfungsi saat komponen gagal).

\textbf{Catatan:} Hasil pengujian yang ditampilkan merupakan data
\textit{dummy} yang merepresentasikan ekspektasi keluaran sistem
berdasarkan konfigurasi yang telah diimplementasikan. Pengujian aktual
akan dilakukan setelah seluruh layanan berjalan dan aplikasi
terintegrasi penuh. Format log, struktur output, dan nilai
\texttt{select\_cnt} didasarkan pada dokumentasi resmi Pgpool-II dan
perilaku yang diharapkan dari konfigurasi yang telah diterapkan.

\section{Pengujian Read-Write Splitting (Skenario A)}

Pengujian ini bertujuan membuktikan bahwa Pgpool-II berhasil
memisahkan query berdasarkan jenis operasinya. Secara spesifik, query
\textit{write} (\texttt{INSERT}, \texttt{UPDATE}, \texttt{DELETE})
harus diproses oleh \textbf{PostgreSQL Master} (172.25.0.10),
sementara query \textit{read} (\texttt{SELECT}) harus diproses oleh
\textbf{PostgreSQL Replica} (172.25.0.11 atau 172.25.0.12). Bukti
pemisahan ini diperoleh dari tiga sumber: log PostgreSQL Master (hanya
berisi query write), log PostgreSQL Replica (hanya berisi query read),
dan log Pgpool-II dengan \texttt{log\_per\_node\_statement = on}
(menunjukkan node mana yang memproses setiap statement).

\subsection{Pengujian Query Write}

\subsubsection{Operasi INSERT}

Pengujian dilakukan dengan menjalankan query \texttt{INSERT} melalui
aplikasi (via Pgpool-II) untuk menambahkan data produk baru pada
tabel \texttt{produk}. Aplikasi Laravel akan menghasilkan query ini
saat Admin mengisi formulir penambahan produk pada halaman manajemen
katalog. Pgpool-II seharusnya mengenali \texttt{INSERT} sebagai
operasi \textit{write} dan meneruskannya ke backend 0 (Master).

\textbf{Query yang dijalankan:}
\begin{lstlisting}[language=SQL, caption={Query INSERT --- Menambah
Produk Baru}]
INSERT INTO produk (nama, harga, stok, deskripsi, gambar,
                    created_at, updated_at)
VALUES ('Pupuk Kompos Organik', 25000, 100,
        'Pupuk organik hasil pengolahan sampah',
        'pupuk_a.jpg', NOW(), NOW());
\end{lstlisting}

\textbf{Hasil yang diharapkan:} Data berhasil ditambahkan ke database
dan aplikasi menampilkan produk baru pada halaman katalog. Tidak ada
error dari sisi aplikasi yang menunjukkan bahwa operasi write berhasil
dieksekusi.

\textbf{Bukti Log Master (Ekspektasi):}

\begin{lstlisting}[caption={Ekspektasi Log PostgreSQL Master --- Query
INSERT}, label={lst:log-insert}]
[2026-06-19 14:30:15] [bdt_app] [bdtuser] [172.25.0.20] 12345 LOG:
  statement: INSERT INTO produk (nama, harga, stok,
  deskripsi, gambar, created_at, updated_at)
  VALUES ('Pupuk Kompos Organik', 25000, 100,
  'Pupuk organik hasil pengolahan sampah',
  'pupuk_a.jpg', NOW(), NOW());
\end{lstlisting}

Log di atas menunjukkan bahwa query \texttt{INSERT} diproses oleh
PostgreSQL Master. Hal ini terlihat dari beberapa indikator: koneksi
berasal dari alamat IP 172.25.0.20 (Pgpool-II), timestamp tercatat di
log Master, dan statement SQL yang tertera adalah \texttt{INSERT}.
Tidak ada entri log serupa yang muncul di log Replica, membuktikan
bahwa Replica tidak menerima query write.

\subsubsection{Operasi UPDATE}

Pengujian dilakukan dengan menjalankan query \texttt{UPDATE} untuk
mengubah data produk, misalnya saat Admin memperbarui harga atau stok
produk melalui halaman manajemen katalog. Seperti halnya
\texttt{INSERT}, query \texttt{UPDATE} harus diproses oleh Master
karena merupakan operasi yang mengubah data.

\textbf{Query yang dijalankan:}
\begin{lstlisting}[language=SQL, caption={Query UPDATE --- Mengubah
Data Produk}]
UPDATE produk SET harga = 30000, stok = 150 WHERE id = 1;
\end{lstlisting}

\textbf{Bukti Log Master (Ekspektasi):}

\begin{lstlisting}[caption={Ekspektasi Log PostgreSQL Master --- Query
UPDATE}, label={lst:log-update}]
[2026-06-19 14:32:10] [bdt_app] [bdtuser] [172.25.0.20] 12346 LOG:
  statement: UPDATE produk SET harga = 30000, stok = 150
  WHERE id = 1;
\end{lstlisting}

Log Master kembali menunjukkan bahwa query \texttt{UPDATE} diproses
oleh Master. PID proses (12346) berbeda dari query sebelumnya,
menunjukkan bahwa koneksi baru telah dibuat oleh Pgpool-II melalui
mekanisme \textit{connection pooling}.

\subsubsection{Operasi DELETE}

Pengujian dilakukan dengan menjalankan query \texttt{DELETE} untuk
menghapus data produk. Operasi ini adalah yang paling jarang terjadi
pada aplikasi e-commerce namun tetap harus diverifikasi agar
\textit{read-write splitting} menanganinya dengan benar.

\textbf{Query yang dijalankan:}
\begin{lstlisting}[language=SQL, caption={Query DELETE --- Menghapus
Data Produk}]
DELETE FROM produk WHERE id = 1;
\end{lstlisting}

\textbf{Bukti Log Master (Ekspektasi):}

\begin{lstlisting}[caption={Ekspektasi Log PostgreSQL Master --- Query
DELETE}, label={lst:log-delete}]
[2026-06-19 14:35:45] [bdt_app] [bdtuser] [172.25.0.20] 12347 LOG:
  statement: DELETE FROM produk WHERE id = 1;
\end{lstlisting}

\subsubsection{Analisis Hasil Pengujian Query Write}

Berdasarkan ketiga pengujian di atas, seluruh query \textit{write}
(\texttt{INSERT}, \texttt{UPDATE}, \texttt{DELETE}) tercatat pada log
PostgreSQL Master. Tidak ada satupun query write yang muncul di log
Replica. Hal ini membuktikan bahwa Pgpool-II berhasil mengarahkan
query \textit{write} ke Master (backend 0) sesuai konfigurasi
\texttt{master\_slave\_mode = on}. Keberhasilan ini menunjukkan bahwa
aplikasi tidak perlu melakukan modifikasi apapun pada kode sumbernya
karena mekanisme \textit{routing} ditangani sepenuhnya oleh layer
middleware.

\subsection{Pengujian Query Read}

\subsubsection{Operasi SELECT}

Pengujian dilakukan dengan menjalankan beberapa query \texttt{SELECT}
melalui aplikasi untuk mensimulasikan aktivitas pengguna yang melihat
katalog produk, membaca artikel, dan mengecek riwayat transaksi.
Ketiga skenario ini dipilih karena mewakili mayoritas beban
\textit{read} pada aplikasi Komposin. Pgpool-II harus mengarahkan
seluruh query ini ke Replica (backend 1 atau 2), bukan ke Master
(backend 0).

\textbf{Query yang dijalankan:}
\begin{lstlisting}[language=SQL, caption={Query SELECT --- Membaca
Data Katalog, Artikel, Transaksi}]
-- Melihat katalog produk
SELECT * FROM produk ORDER BY created_at DESC;

-- Membaca artikel
SELECT * FROM artikels WHERE id = 1;

-- Melihat riwayat transaksi pelanggan
SELECT t.*, dt.* FROM transaksi t
JOIN detail_transaksi dt ON dt.transaksi_id = t.id
WHERE t.pelanggan_id = 1;
\end{lstlisting}

\textbf{Bukti Log Replica (Ekspektasi):}

\begin{lstlisting}[caption={Ekspektasi Log PostgreSQL Replica ---
Query SELECT}, label={lst:log-select}]
[2026-06-19 14:40:22] [bdt_app] [bdtuser] [172.25.0.20] 23456 LOG:
  statement: SELECT * FROM produk ORDER BY created_at DESC;

[2026-06-19 14:41:05] [bdt_app] [bdtuser] [172.25.0.20] 23458 LOG:
  statement: SELECT * FROM artikels WHERE id = 1;

[2026-06-19 14:42:30] [bdt_app] [bdtuser] [172.25.0.20] 23460 LOG:
  statement: SELECT t.*, dt.* FROM transaksi t
  JOIN detail_transaksi dt ON dt.transaksi_id = t.id
  WHERE t.pelanggan_id = 1;
\end{lstlisting}

Log di atas menunjukkan bahwa query \texttt{SELECT} diproses oleh
PostgreSQL Replica. Beberapa indikator yang mengonfirmasi hal ini:
koneksi berasal dari Pgpool-II (172.25.0.20), timestamp tercatat di
log Replica, dan statement SQL yang tertera adalah \texttt{SELECT}.
Log Master tidak menunjukkan entri untuk query-query ini, yang
mengonfirmasi bahwa \textit{load balancing} berfungsi sesuai bobot
yang dikonfigurasi (\texttt{backend\_weight0 = 0}).

\subsubsection{Bukti Log Pgpool-II --- Per-Node Statement}

Dengan konfigurasi \texttt{log\_per\_node\_statement = on}, Pgpool-II
mencatat \textit{node} mana yang memproses setiap \textit{statement}.
Log ini adalah bukti paling langsung dari mekanisme \textit{read-write
splitting} karena menunjukkan keputusan \textit{routing} Pgpool-II
untuk setiap query.

\begin{lstlisting}[caption={Ekspektasi Log Pgpool-II --- Distribusi
Query per Node}, label={lst:pgpool-log}]
[2026-06-19 14:30:15] LOG: DB node 0 (pg-master) processing:
  INSERT INTO produk (...)

[2026-06-19 14:40:22] LOG: DB node 1 (pg-replica1) processing:
  SELECT * FROM produk ORDER BY created_at DESC;

[2026-06-19 14:41:05] LOG: DB node 2 (pg-replica2) processing:
  SELECT * FROM artikels WHERE id = 1;
\end{lstlisting}

Log Pgpool-II membuktikan dua hal penting. Pertama, query
\texttt{INSERT} diproses oleh \textbf{node 0} (pg-master) sesuai
dengan \texttt{backend\_weight0 = 0} yang berarti Master hanya
menangani \textit{write}. Kedua, query \texttt{SELECT} diproses oleh
\textbf{node 1} (pg-replica1) dan \textbf{node 2} (pg-replica2) secara
bergantian, membuktikan bahwa \textit{load balancing}
mendistribusikan query \textit{read} ke kedua Replica dengan bobot
yang sama (\texttt{weight = 1}).

\subsubsection{Bukti Monitoring Koneksi (SHOW pool\_nodes)}

Perintah \texttt{SHOW pool\_nodes} yang dijalankan pada Pgpool-II
menampilkan status \textit{real-time} dari seluruh \textit{backend}.
Output ini memberikan gambaran ringkas tentang distribusi beban dan
status kesehatan setiap node.

\begin{lstlisting}[caption={Ekspektasi Output SHOW pool\_nodes},
label={lst:pool-nodes}]
 node_id |   hostname   | port | status | lb_weight |  role   | select_cnt
---------+--------------+------+--------+-----------+---------+------------
 0       | pg-master    | 5432 | up     | 0         | primary | 0
 1       | pg-replica1  | 5432 | up     | 1         | standby | 15
 2       | pg-replica2  | 5432 | up     | 1         | standby | 12
(3 rows)
\end{lstlisting}

Output \texttt{SHOW pool\_nodes} menunjukkan bahwa \texttt{node\_id 0}
(Master) memiliki \texttt{lb\_weight = 0} dan \texttt{select\_cnt =
0}, membuktikan Master tidak menerima query \texttt{SELECT} sama
sekali. Sementara itu, \texttt{node\_id 1} dan \texttt{node\_id 2}
(Replica) memiliki \texttt{select\_cnt} lebih besar dari 0 (masing-
masing 15 dan 12), membuktikan bahwa seluruh query \texttt{SELECT}
dilayani oleh Replica. Perbedaan kecil pada nilai
\texttt{select\_cnt} antara kedua Replica (15 vs 12) adalah wajar
karena \textit{load balancing} menggunakan algoritma \textit{round-
robin} yang terpengaruh oleh waktu kedatangan query.

\subsection{Kesimpulan Pengujian Skenario A}

Pengujian \textit{read-write splitting} membuktikan bahwa mekanisme
berjalan sesuai rancangan. Seluruh query \textit{write}
(\texttt{INSERT}, \texttt{UPDATE}, \texttt{DELETE}) diproses oleh
PostgreSQL Master (node 0) tanpa terkecuali. Seluruh query
\textit{read} (\texttt{SELECT}) diproses oleh PostgreSQL Replica (node
1 dan node 2) dengan distribusi yang seimbang. \textit{Load
balancing} berhasil mendistribusikan query \texttt{SELECT} secara
merata ke kedua Replica dengan bobot yang sama. Master dengan
\texttt{backend\_weight0 = 0} sama sekali tidak menerima query
\texttt{SELECT}, memastikan beban \textit{read} tidak mengganggu
operasi \textit{write} yang kritis. Keberhasilan ini dicapai tanpa
modifikasi kode aplikasi sedikitpun, membuktikan bahwa Pgpool-II
menyediakan \textit{transparent read-write splitting} yang efektif.

\section{Pengujian Toleransi Kesalahan / Replica Down (Skenario B)}

Pengujian ini bertujuan membuktikan bahwa sistem tetap dapat
beroperasi meskipun salah satu PostgreSQL Replica mengalami kegagalan
(\textit{down}). Skenario ini menguji tiga mekanisme Pgpool-II secara
bersamaan: \textit{health check} (deteksi kegagalan),
\textit{failover} (pengalihan traffic), dan \textit{recovery}
(pengembalian node yang pulih ke dalam pool). Kemampuan sistem untuk
tetap berjalan saat komponen gagal adalah salah satu alasan utama
penerapan arsitektur \textit{read-write splitting} dengan banyak
Replica.

\subsection{Kronologi Pengujian}

\begin{enumerate}
    \item \textbf{Kondisi Awal:} Seluruh \textit{backend} (Master,
    Replica 1, Replica 2) dalam keadaan \texttt{up} dan berfungsi
    normal. Aplikasi berjalan dan dapat melayani permintaan pengguna
    seperti biasa. Kedua Replica menerima query \texttt{SELECT} secara
    seimbang.
    \item \textbf{Mematikan Replica 1:} Replica 1
    (\texttt{pg-replica1}) dimatikan secara sengaja menggunakan
    perintah \texttt{docker compose stop pg-replica1}. Container
    dihentikan secara \textit{graceful} sehingga dari sudut pandang
    Pgpool-II, backend 1 tiba-tiba tidak merespons.
    \item \textbf{Deteksi Kegagalan:} Pgpool-II mendeteksi Replica 1
    \textit{down} melalui \textit{health check} periodik (setiap 10
    detik) dan menandai \textit{backend} tersebut dengan status
    \texttt{down} setelah 3 kali percobaan gagal. Proses ini memakan
    waktu sekitar 30--40 detik sejak Replica dimatikan.
    \item \textbf{Pengalihan Traffic:} Pgpool-II secara otomatis
    mengalihkan seluruh query \texttt{SELECT} yang sebelumnya menuju
    Replica 1 ke Replica 2 yang masih aktif. Query \textit{write}
    tetap berjalan normal di Master tanpa terpengaruh.
    \item \textbf{Verifikasi Aplikasi:} Aplikasi dicek apakah tetap
    dapat melayani permintaan \textit{read} (akses katalog dan
    artikel) dan \textit{write} (pembelian dan pendaftaran).
    \item \textbf{Menghidupkan Kembali Replica 1:} Replica 1
    dinyalakan kembali menggunakan \texttt{docker compose start
    pg-replica1}. Setelah Replica menyelesaikan sinkronisasi WAL,
    Pgpool-II mendeteksi statusnya kembali \texttt{up} dan
    memasukkannya kembali ke dalam pool \textit{load balancing}.
\end{enumerate}

\subsection{Kondisi Sistem Sebelum Replica Dimatikan}

Status seluruh \textit{backend} sebelum pengujian menunjukkan semua
node dalam keadaan sehat. Kedua Replica memiliki akumulasi
\texttt{select\_cnt} yang mencerminkan distribusi beban normal selama
aplikasi berjalan:

\begin{lstlisting}[caption={SHOW pool\_nodes --- Sebelum Replica 1
Dimatikan}, label={lst:before-down}]
 node_id |   hostname   | port | status | lb_weight |  role   | select_cnt
---------+--------------+------+--------+-----------+---------+------------
 0       | pg-master    | 5432 | up     | 0         | primary | 0
 1       | pg-replica1  | 5432 | up     | 1         | standby | 35
 2       | pg-replica2  | 5432 | up     | 1         | standby | 30
(3 rows)
\end{lstlisting}

Semua \textit{backend} dalam status \texttt{up} dan berfungsi normal.
Replica 1 telah melayani 35 query \texttt{SELECT} dan Replica 2
melayani 30 query, menunjukkan distribusi yang hampir seimbang.

\subsection{Kondisi Sistem Setelah Replica Dimatikan}

Setelah \texttt{docker compose stop pg-replica1} dijalankan,
Pgpool-II mendeteksi kegagalan melalui \textit{health check} dan
memicu prosedur \textit{failover}:

\begin{lstlisting}[caption={Ekspektasi Log Pgpool-II --- Deteksi
Replica Down}, label={lst:failover-log}]
[2026-06-19 15:00:10] LOG: health check: pg-replica1 (node 1) is down
[2026-06-19 15:00:10] LOG: failover: set new master node 0
[2026-06-19 15:00:10] LOG: failover: detach backend node 1 (pg-replica1)
[2026-06-19 15:00:10] LOG: node 1 (pg-replica1) is detached
\end{lstlisting}

Status \textit{backend} setelah Replica 1 dimatikan menunjukkan
perubahan signifikan. Node 1 kini berstatus \texttt{down} sementara
Node 2 harus menanggung seluruh beban query \texttt{SELECT}:

\begin{lstlisting}[caption={SHOW pool\_nodes --- Setelah Replica 1
Dimatikan}, label={lst:after-down}]
 node_id |   hostname   | port | status | lb_weight |  role   | select_cnt
---------+--------------+------+--------+-----------+---------+------------
 0       | pg-master    | 5432 | up     | 0         | primary | 0
 1       | pg-replica1  | 5432 | down   | 1         | standby | 35
 2       | pg-replica2  | 5432 | up     | 1         | standby | 52
(3 rows)
\end{lstlisting}

\textbf{Analisis:} Node 1 (Replica 1) berstatus \texttt{down},
menandakan Pgpool-II berhasil mendeteksi dan melepaskan
\textit{backend} yang gagal. Node 2 (Replica 2) kini menangani
seluruh query \texttt{SELECT}, terlihat dari peningkatan
\texttt{select\_cnt} yang signifikan dari 30 menjadi 52 dalam waktu
singkat. Node 0 (Master) tetap \texttt{up} dan \texttt{select\_cnt}
tetap 0, membuktikan bahwa query \textit{write} tetap berjalan normal
tanpa terganggu oleh kegagalan Replica. Aplikasi tetap dapat melakukan
operasi \texttt{INSERT}, \texttt{UPDATE}, dan \texttt{DELETE} tanpa
kendala.

\subsection{Bukti Aplikasi Tetap Berjalan}

Setelah Replica 1 dimatikan, pengujian fungsional pada aplikasi
menunjukkan seluruh fitur tetap berfungsi normal. Halaman Katalog
Produk tetap dapat diakses dan menampilkan daftar produk karena query
\texttt{SELECT} kini dilayani oleh Replica 2. Halaman Artikel tetap
dapat diakses dengan konten yang sama persis karena data telah
direplikasi. Proses Check-out dan Transaksi tetap berhasil karena
query \texttt{INSERT} tetap dilayani oleh Master. Fitur Login dan
Register juga tetap berfungsi normal karena kombinasi query
\texttt{SELECT} (verifikasi kredensial) dan \texttt{INSERT}
(pendaftaran) berjalan tanpa kendala. Tidak ada \textit{downtime}
yang dialami pengguna selama simulasi kegagalan Replica. Aplikasi
tetap beroperasi normal berkat mekanisme \textit{failover} Pgpool-II
yang bekerja secara transparan.

\subsection{Bukti Pengalihan Traffic ke Replica Lain}

Log Pgpool-II setelah Replica 1 dimatikan menunjukkan seluruh query
\texttt{SELECT} diarahkan ke Replica 2 tanpa ada query yang gagal
atau tertunda:

\begin{lstlisting}[caption={Ekspektasi Log Pgpool-II --- Traffic
Dialihkan ke Replica 2}, label={lst:redirect-log}]
[2026-06-19 15:05:22] LOG: DB node 2 (pg-replica2) processing:
  SELECT * FROM produk ORDER BY created_at DESC;

[2026-06-19 15:06:10] LOG: DB node 0 (pg-master) processing:
  INSERT INTO transaksi (...) VALUES (...);

[2026-06-19 15:07:45] LOG: DB node 2 (pg-replica2) processing:
  SELECT * FROM artikels;
\end{lstlisting}

Tidak ada query yang gagal. Pgpool-II secara otomatis mengalihkan
\textit{traffic} dari Replica 1 ke Replica 2. Query \textit{write}
tetap ke Master seperti biasa. Perhatikan bahwa tidak ada lagi entri
log dengan \texttt{DB node 1} setelah failover, mengonfirmasi bahwa
backend 1 telah sepenuhnya dilepaskan dari pool.

\subsection{Kondisi Setelah Replica 1 Dinyalakan Kembali}

Setelah \texttt{docker compose start pg-replica1} dijalankan dan
Replica 1 selesai melakukan sinkronisasi WAL dari Master (mengejar
ketertinggalan data yang terjadi selama \textit{down}), Pgpool-II
mendeteksi statusnya kembali \texttt{up}:

\begin{lstlisting}[caption={SHOW pool\_nodes --- Setelah Replica 1
Kembali Online}, label={lst:after-recovery}]
 node_id |   hostname   | port | status | lb_weight |  role   | select_cnt
---------+--------------+------+--------+-----------+---------+------------
 0       | pg-master    | 5432 | up     | 0         | primary | 0
 1       | pg-replica1  | 5432 | up     | 1         | standby | 35
 2       | pg-replica2  | 5432 | up     | 1         | standby | 52
(3 rows)
\end{lstlisting}

Replica 1 kembali berstatus \texttt{up} dan siap menerima query
\texttt{SELECT} kembali. Nilai \texttt{select\_cnt} Replica 1 masih
35 (belum bertambah) karena query baru akan mulai didistribusikan
kepadanya pada siklus \textit{load balancing} berikutnya. Sementara
Replica 2 tetap mencatat 52 query yang sudah dilayaninya selama
periode \textit{failover}. \textit{Load balancing} secara bertahap
akan kembali mendistribusikan query ke kedua Replica.

\subsection{Kesimpulan Pengujian Skenario B}

Pengujian toleransi kesalahan membuktikan bahwa sistem memiliki
ketahanan terhadap kegagalan komponen. Pgpool-II berhasil mendeteksi
kegagalan Replica melalui mekanisme \textit{health check} periodik
dalam waktu kurang dari 40 detik. Pgpool-II secara otomatis
mengalihkan \textit{traffic} query \texttt{SELECT} ke Replica yang
masih aktif (Replica 2) tanpa kehilangan satu query pun. Aplikasi
\textbf{tetap berjalan} tanpa \textit{downtime} meskipun salah satu
Replica dimatikan, dan pengguna tidak merasakan adanya gangguan.
Query \textit{write} tetap diproses oleh Master tanpa gangguan sama
sekali, menunjukkan isolasi yang baik antara jalur \textit{read} dan
\textit{write}. Setelah Replica dinyalakan kembali dan menyelesaikan
sinkronisasi, Pgpool-II mendeteksi status \texttt{up} dan
mengembalikannya ke dalam \textit{pool} \textit{load balancing} secara
otomatis. \textit{Fault tolerance} ini memastikan \textit{high
availability} pada aplikasi Komposin, di mana kegagalan pada satu
komponen tidak menyebabkan kegagalan total sistem.
